\documentclass[11pt,a4paper]{article}
\usepackage[utf8]{inputenc}\usepackage[T1]{fontenc}
\usepackage[provide=*,english]{babel}
\usepackage{lmodern,microtype,amsmath,amssymb,booktabs,geometry,xcolor,fancyhdr}
\usepackage[hidelinks]{hyperref}
\geometry{margin=2.3cm,headheight=14pt}
\setlength{\parindent}{0pt}\setlength{\parskip}{0.4em}
\definecolor{gcblue}{RGB}{34,73,110}\definecolor{gcgray}{RGB}{90,96,102}
\pagestyle{fancy}\fancyhf{}
\fancyhead[L]{\small\color{gcgray}GC2D: HBVM(4,2)}
\fancyhead[R]{\small\color{gcgray}Model theory}\fancyfoot[C]{\thepage}
\setlength{\emergencystretch}{3em}
\begin{document}
\begin{center}
 {\LARGE\bfseries\color{gcblue}Hamiltonian Boundary Value Method HBVM(4,2)}\par
 \vspace{0.35em}{\large Rank-two polynomial path with four Gauss nodes}
\end{center}

\section{Polynomial-path formulation}

Let $(c_i,b_i)_{i=1}^4$ be four-point Gauss--Legendre quadrature on $[0,1]$.
For the first two orthonormal shifted Legendre polynomials,
\[
P_0(c)=1,\qquad P_1(c)=\sqrt3(2c-1),
\]
define
\[
I_{i0}=c_i,\qquad I_{i1}=\sqrt3(c_i^2-c_i).
\]
HBVM(4,2) solves for two vector coefficients $\gamma_0,\gamma_1$:
\begin{align}
Y_i&=y_n+h\sum_{j=0}^1I_{ij}\gamma_j,\\
R_j(\gamma)&=\gamma_j-
\sum_{i=1}^4b_iP_j(c_i)f(t_n+c_ih,Y_i)=0.
\end{align}
The accepted state is
\[
y_{n+1}=y_n+h\sum_{i=1}^4b_i f(t_n+c_ih,Y_i).
\]
Equivalently, the Runge--Kutta matrix is
$A=\mathcal I_2\mathcal P_2^T\Omega_q$ and has rank two.  The nonlinear
unknown therefore has twice the state dimension, although four stages evaluate
the line integral.

\section{Order and energy property}

With polynomial degree $s=2$, the method has classical order $2s=4$.
Four-point Gauss quadrature has degree seven.  For an autonomous polynomial
Hamiltonian of degree $\nu$, exact polynomial energy conservation follows when
$k\geq\nu s/2$; hence HBVM(4,2) covers degrees through four.  For a general
smooth Hamiltonian the quadrature error makes energy conservation approximate,
with an error that decreases rapidly with the step and quadrature accuracy.

HBVM$(k,s)$ is generally not symplectic when $k>s$.  Thus HBVM(4,2) must not be
labeled symplectic by construction; map-Jacobian diagnostics measure its
canonical defect independently of its energy behavior.

\section{Nonlinear configurations}

Newton acts on the reduced residual $(R_0,R_1)$.  The public
\texttt{jacobian\_method} values are \texttt{analytic},
\texttt{finite\_difference}, and \texttt{auto}.  Analytic mode assembles exact
guiding-centre particle blocks at the four stages; finite-difference mode
differentiates the stage fields; auto selects the available capability.  A
backtracking correction is used only when the full Newton update increases the
residual.  The stopping scale is
\[
\varepsilon_{\rm abs}+\varepsilon_{\rm rel}
\max(1,\lVert y_n\rVert_\infty).
\]

With \texttt{track\_energy=True}, the implementation stores the physical
Hamiltonian and its drift at saved times.  For explicitly time-dependent
systems this is a diagnostic, not an invariant claim.

\section{Implementation correspondence}

Coefficients, residuals, Newton assembly, and integration are in
\path{src/methods/hbvm/order4.py}.  The simulation document describes
required Hamiltonian capabilities, work counters, fixed-grid sampling, and
evaluation studies; shared physical dynamics are documented outside the model
tree.

\section{Four state formulations and passive energy tracking}
For one planar particle the supported internal representations are
\[
\begin{array}{lll}
\text{physical} & z=(x,y) & \mathbb{R}^{2},\\
\text{physical with energy} & (z,t,\kappa) & \mathbb{R}^{4},\\
\text{duplicated} & (u,v),\quad u,v\in\mathbb{R}^{2} & \mathbb{R}^{4},\\
\text{duplicated with energy} & (u,v,t,\kappa) & \mathbb{R}^{6}.
\end{array}
\]
Classical methods use the physical rows; ABBA and BM4 use the duplicated rows.
\texttt{track\_energy} selects the energy variant. For $N$ planar particles the
respective dimensions are $2N$, $4N$, $4N$, and $6N$: one time entry and
one momentum per particle. All time entries equal the integration time. Classical full-cyclotron runs instead use $4N$ or $6N$ entries.
Accepted spatial copies coincide; only one physical copy is exported.
Every coordinate block contains $N$ values in component-major order. The
formulation validates particle times and aligns output round-off; the
integrator does not interpret coordinate indices.

The normalized passive balance is
\[
\dot\kappa=-\partial_tH(t,z),\qquad \kappa(t_0)=0,\qquad
\Delta K(t)=H(t,z(t))+\kappa(t)-H(t_0,z_0).
\]
Splitting accumulators from both copies are divided by two. The physical
Hamiltonian may vary with time; $\Delta K$ measures its computed balance.
Tracking leaves physical trajectories, nonlinear stopping rules and adaptive
acceptance unchanged. Dynamics must supply the momentum derivative, including
zero for an autonomous Hamiltonian.

Hairer's constraint is solely $u-v=0$: the reduced multiplier has $2N$ entries
and ABBA's simultaneous spatial solve has $6N$ unknowns. Neither $t$ nor
$\kappa$ is projected. The former eight-component projection is retired;
\texttt{state\_extension='fully\_extended'} raises migration guidance.

Dynamics, formulation, method and integration remain separate. The shared
state classes own layouts, diagonal embedding and history extraction in
\begin{quote}\small\path{src/formulations/state.py}\end{quote}
Each integration initializes a fresh method instance. Observers receive physical
states; energy histories are in \texttt{Solution.diagnostics}.
The \texttt{extended\_time} history has shape $(N,n_{\rm samples})$. Fixed-step shadow
samples leave main states and counters unchanged. Adaptive energy quadrature
uses accepted dense output; its backend state and Jacobian remain physical.

The four converged stage states and Gauss weights integrate the passive momentum. Its value never enters the Legendre coefficient solve.

\end{document}
